La organización del proyecto se ha definido atendiendo a su naturaleza académica y al hecho de que el desarrollo ha sido realizado por un único estudiante. A pesar de ello, el proyecto cuenta con distintos perfiles implicados, tanto desde el punto de vista de supervisión académica como desde el conocimiento del dominio del problema.

\subsection{Personas y roles involucrados}

Los principales roles involucrados en el proyecto son los siguientes:

\begin{itemize}
    \item \textbf{Autor del proyecto:} responsable del análisis, diseño, implementación, pruebas y documentación del sistema. Asume las tareas propias de analista, diseñador, desarrollador y responsable de documentación.
    \item \textbf{Director del Trabajo Fin de Grado:} encargado de orientar el desarrollo académico del proyecto, revisar el avance de la memoria y supervisar que el trabajo se ajuste a los objetivos propios de un Trabajo Fin de Grado.
    \item \textbf{Cliente o usuario experto:} distribuidor de productos capilares profesionales que aporta información sobre la situación actual, los procesos de negocio y las necesidades reales del dominio.
    \item \textbf{Usuarios finales previstos:} perfiles que podrían utilizar el sistema una vez evolucionado, incluyendo administrador o distribuidor, comerciales y peluquerías profesionales.
\end{itemize}

La comunicación con el cliente o usuario experto ha sido especialmente relevante durante la fase de captura de requisitos, ya que ha permitido comprender la operativa real del distribuidor y detectar las carencias de los procesos actuales. Por su parte, las revisiones académicas han servido para orientar la estructura de la memoria y ajustar el alcance del sistema al contexto del proyecto.

\subsection{Estructura de trabajo}

Al tratarse de un proyecto individual, no se ha definido una estructura jerárquica compleja. Sin embargo, las responsabilidades se han organizado por áreas de trabajo:

\begin{itemize}
    \item \textbf{Gestión del proyecto:} planificación temporal, control del alcance, seguimiento de hitos y gestión de riesgos.
    \item \textbf{Análisis funcional:} estudio del dominio, identificación de actores, definición de requisitos y reglas de negocio.
    \item \textbf{Diseño técnico:} definición de arquitectura, diseño de datos, diseño de interfaz y selección tecnológica.
    \item \textbf{Construcción:} implementación de la aplicación, configuración del entorno, base de datos y funcionalidades iniciales.
    \item \textbf{Validación:} realización de pruebas, revisión de flujos y comprobación del cumplimiento de los requisitos desarrollados.
    \item \textbf{Documentación:} elaboración de la memoria, anexos, manuales y documentación técnica asociada.
\end{itemize}

\subsection{Recursos utilizados}

Para el desarrollo del proyecto se han empleado recursos hardware, software y servicios externos. La Tabla~\ref{tab:project-resources} resume los principales recursos utilizados.

\begin{table}[htbp]
    \centering
    \small
    \caption{Recursos utilizados en el proyecto}
    \label{tab:project-resources}
    \begin{tabularx}{\textwidth}{|L{3.2cm}|L{3.5cm}|Y|}
        \hline
        \textbf{Tipo de recurso} & \textbf{Recurso} & \textbf{Uso en el proyecto} \\
        \hline
        Hardware & Equipo de desarrollo & Implementación, pruebas locales, documentación y ejecución del entorno de desarrollo. \\
        \hline
        Software base & Windows 11 & Sistema operativo utilizado para el desarrollo del proyecto. \\
        \hline
        Entorno de desarrollo & Visual Studio Code & Edición de código fuente, documentación LaTeX y gestión del proyecto. \\
        \hline
        Control de versiones & Git & Registro de cambios, seguimiento de versiones y control del código fuente. \\
        \hline
        Repositorio remoto & GitHub & Almacenamiento remoto del código fuente y respaldo del proyecto. \\
        \hline
        Base de datos & PostgreSQL & Persistencia de los datos principales de la aplicación. \\
        \hline
        Contenedores & Docker y Docker Compose & Ejecución de servicios locales y reproducción del entorno de desarrollo. \\
        \hline
        Documentación & LaTeX & Elaboración de la memoria académica del Trabajo Fin de Grado. \\
        \hline
        Diagramas & Draw.io e Inkscape & Elaboración y adaptación de diagramas incluidos en la documentación. \\
        \hline
        Pruebas externas & Port Forwarding de Visual Studio Code (Microsoft Dev Tunnels) & Acceso temporal al entorno local desde otros dispositivos durante las pruebas, con una URL publica reutilizable comoda para validar la PWA en móvil. \\
        \hline
    \end{tabularx}
\end{table}

La combinación de estos recursos ha permitido desarrollar una solución reproducible, documentada y preparada para continuar su evolución en fases posteriores.
